\documentclass[11pt]{article}
\usepackage[a4paper,margin=23mm]{geometry}
\usepackage{amsmath,amssymb,amsthm,mathtools,bm}
\usepackage{booktabs,longtable,array}
\usepackage{enumitem}
\usepackage{xcolor,microtype}
\usepackage[hidelinks]{hyperref}
\setlength{\emergencystretch}{3em}
\newtheorem{theorem}{Theorem}[section]
\newtheorem{proposition}[theorem]{Proposition}
\newtheorem{lemma}[theorem]{Lemma}
\theoremstyle{definition}
\newtheorem{gate}[theorem]{Fail-closed gate}
\theoremstyle{remark}
\newtheorem{remark}[theorem]{Remark}
\newcommand{\RR}{\mathbb R}
\newcommand{\CC}{\mathbb C}
\newcommand{\one}{\mathbf 1}
\newcommand{\dd}{\mathrm d}
\newcommand{\sgn}{\operatorname{sgn}}
\newcommand{\diag}{\operatorname{diag}}
\newcommand{\cC}{\mathcal C}
\newcommand{\cH}{\mathcal H}

\title{AP-E7 Discrete Re-relaxation Audit:\\
Finite-site Topology, an Admissibility Obstruction,\\
and Same-grid Sparse Operators}
\author{Route F research note}
\date{17 July 2026}

\begin{document}
\maketitle

\begin{abstract}
The AP-E6 continuum hedgehog is used only as an initial condition. We then
minimize the same declared finite-difference $n+s$ energy independently on
five cubic grids with exact-vacuum Dirichlet data. This exposes a structural
obstruction: at finite site number the unconstrained configuration space is
path-connected, so it has no exact $B\in\pi_3(S^3)$ components. A
Freudenthal-tetrahedron regular-value degree equals $B=1$ only on an
admissible open subset. Unconstrained solves cross its dislocation boundary
and converge to $B=0$. Endpoint-guarded line searches at three explicit
admissibility floors have $B=1$ at every accepted iterate, but hit the floor
with nonzero projected gradient instead of finding an interior stationary
point. Continuous admissibility of the entire geodesic segment between two
accepted endpoints is not certified.

For numerical and algebraic controls only, a central $3^3$ hedgehog core is
fixed and every other variable is re-relaxed. In that artificial sector we
assemble an exact sparse constraint-projected $n+s$ Hessian and same-grid
diquark, gauge, ghost, and Wilson/Yukawa normal operators. The latter is a
positive fermion proxy, not the bosonic second derivative of
$-\log\det D_W$. Therefore no genuine unanchored relaxed $B=1$ lattice
solution, physical super-Hessian, four-dimensional dynamics, or continuum
stability theorem is claimed.
\end{abstract}

\tableofcontents

\section{Declared finite-difference problem}

Let $n_x\in S^3\subset\RR^4$, $s_x\in\RR$, and freeze every boundary site to
$n_x=(1,0,0,0)$, $s_x=1$. The AP-E6 energy is discretized as
\begin{align}
 E_a[n,s]={}&\frac a2\sum_{\langle xy\rangle}
 \left(\lVert n_y-n_x\rVert^2+(s_y-s_x)^2\right)\nonumber\\
 &+\frac{\kappa a^3}{2}\sum_c\sum_{i<j}
 \left(\lVert u_i^c\rVert^2\lVert u_j^c\rVert^2
 -(u_i^c\!\cdot u_j^c)^2\right)\nonumber\\
 &+a^3\sum_{x\in\mathrm{int}}
 \left[\frac{m_s^2}{2}(s_x-1)^2
 +\beta^2s_x^2(1-n_{x0})\right],
 \label{eq:lattice-energy}
\end{align}
where $u_i^c=(n_{c+\hat\imath}-n_c)/a$ and
\begin{equation}
 \kappa=1,\qquad \beta=\frac12,\qquad m_s=\frac52.
 \label{eq:parameters}
\end{equation}
This is an ordinary classical three-dimensional regulator of the continuum
Skyrme-type model \cite{Skyrme1961APE7}. It is neither a Euclidean
four-dimensional measure nor a two-colour lattice-QCD action. Recent
two-colour sigma-model studies motivate additional scalar channels but do not
remove this distinction \cite{SuenagaEtAl2023APE7,SuenagaEtAl2024APE7}.

Each grid is initialized by the same AP-E6 $R=16$ continuum solution with
canonical little-endian float64 checksum
\begin{center}
\path{81104f59a5f3f4337739fc2a217cafc8da91581c9f1fb40b4fdba6ce0f948d59}.
\end{center}
The sampled outer faces are replaced by the exact vacuum before each
independent lattice solve.

\section{Why a finite-site \texorpdfstring{$B$}{B} sector is not automatic}

\begin{theorem}[Finite-site no-sector theorem]
For $M=(N-2)^3<\infty$ dynamical sites and fixed boundary, the unconstrained
configuration space
\begin{equation}
 \cC_N=(S^3)^M\times\RR^M
 \label{eq:configuration-space}
\end{equation}
is path-connected. Consequently no integer-valued continuous function on all
of $\cC_N$ can distinguish a $B=1$ component from the vacuum.
\end{theorem}

\begin{proof}
Both $S^3$ and $\RR$ are path-connected. Given two configurations, choose a
path between the two values at every site and take their finite Cartesian
product. This constructs a path in $\cC_N$. A continuous map from a connected
space to the discrete space $\mathbb Z$ is constant. Thus continuum
$\pi_3(S^3)$ cannot become an exact decomposition of the unrestricted
finite-site space without an additional admissibility condition,
topology-preserving interpolation, or infinite dislocation barrier.
\end{proof}

This is the finite-site form of the dislocation problem familiar from
geometric lattice charge \cite{BergLuscher1981APE7,Berg1981DislocationsAPE7}.
Lattice Skyrme constructions that retain metastable topology introduce an
explicit spherical-tetrahedron prescription and stabilizing structure
\cite{SchrammSvetitsky2000APE7}.

\section{Freudenthal regular-value degree}

Split every cube into the six Freudenthal tetrahedra indexed by permutations
$p\in S_3$, with vertices
\begin{equation}
 x_0=c,\quad x_1=c+e_{p_1},\quad
 x_2=c+e_{p_1}+e_{p_2},\quad x_3=c+(1,1,1).
\end{equation}
Their domain orientation is $\epsilon(p)=\sgn(p)$. For a generic regular
target $q\in S^3$, form
\begin{equation}
 V_T=[n_{x_0}\ n_{x_1}\ n_{x_2}\ n_{x_3}],\qquad
 c_T=V_T^{-1}q\quad\text{when }\det V_T\ne0.
 \label{eq:target-matrix}
\end{equation}
There are finitely many tetrahedra.  If $V_T$ is singular, its column span is
a proper linear subspace of $\mathbb R^4$; choose the generic regular target
$q$ outside the finite union of those spans.  Such a tetrahedron then has no
preimage of $q$ and contributes zero.  Equivalently, an exceptional target is
defined by an arbitrarily small regular perturbation before taking the signed
count.  Thus the inverse in \eqref{eq:target-matrix} is never applied to a
singular matrix.  The numerical implementation uses the declared
$|\det V_T|<10^{-13}$ exceptional-cell tolerance and verifies target agreement
and a much larger nonzero determinant margin on every counted preimage.
The ray through $q$ intersects the positive cone of the four vertices iff
every component of $c_T$ is positive. In that case the local map degree is
\begin{equation}
 d_T(q)=\epsilon(p)\sgn\det V_T.
\end{equation}
Thus, in the sign convention of the continuum AP-E6 card,
\begin{equation}
 \deg(n;q)=\sum_T d_T(q),\qquad B_{\rm geom}=-\deg(n;q).
 \label{eq:degree}
\end{equation}
Three generic targets are used, and agreement is required.

The normalized affine map on $T$ is defined only if the vertex convex hull
does not contain the origin. Its exact margin is
\begin{equation}
 \delta_T=\min_{\lambda_i\ge0,\ \sum_i\lambda_i=1}
 \left\lVert\sum_{i=0}^3\lambda_i n_{x_i}\right\rVert,
 \qquad \delta_a=\min_T\delta_T.
 \label{eq:admissibility}
\end{equation}
The implementation enumerates all active faces of each four-vertex convex
program. When $\delta_a>0$, normalization is nonsingular and
$B_{\rm geom}$ is locally constant. At $\delta_a=0$, a dislocation can change
the integer without contradicting the theorem above. A regression uses the
four regular-tetrahedron vectors
$(1,1,1,0)$, $(1,-1,-1,0)$, $(-1,1,-1,0)$, and $(-1,-1,1,0)$, divided by
$\sqrt3$. Their equal-weight sum vanishes; the active-set KKT solver returns
$\delta_a=2.21\times10^{-16}$ and correctly rejects admissibility.

\section{Riemannian relaxation and its audit}

At every site choose an orthonormal tangent frame $E_x\in\RR^{4\times3}$.
A fixed-base exponential chart is
\begin{equation}
 n_x(z_x)=\cos r_x\,n_x^{(0)}
 +\frac{\sin r_x}{r_x}E_xz_x,
 \qquad r_x=\lVert z_x\rVert,
 \label{eq:chart}
\end{equation}
while $s_x=s_x^{(0)}+\sigma_x$. If $u=z/r$, the exact pullback of an
embedding gradient $g\in\RR^4$ is
\begin{align}
 J(z)^Tg={}&u\bigl[(-\sin r\,n^{(0)}+\cos r\,Eu)\cdot g\bigr]\nonumber\\
 &+\frac{\sin r}{r}(I-uu^T)E^Tg,
 \label{eq:chart-gradient}
\end{align}
with limit $E^Tg$ at $r=0$. A centered directional derivative checks
Eq.~\eqref{eq:chart-gradient} independently; its best relative residual is
$7.31\times10^{-11}$.

Three logically different solves are reported:
\begin{enumerate}[label=(\alph*)]
\item unconstrained L-BFGS relaxation of every interior variable;
\item Riemannian Armijo descent that rejects any candidate endpoint with
  $B_{\rm geom}\ne1$ or $\delta_a<\delta_{\rm floor}$;
\item an artificial control fixing the $n$ values, but not $s$, on the
  central $3\times3\times3$ sampled hedgehog block.
\end{enumerate}
Only (a) tests full-grid stationarity. Case (b) diagnoses the admissible
boundary. Case (c) tests a deliberately restricted solver and cannot be
called a genuine relaxed lattice soliton.

\section{Exact restricted Hessian and same-grid blocks}

Let $G_x=\partial E_a/\partial n_x$ and $\lambda_x=n_x\cdot G_x$. With
$Q=\diag(E_1,1,E_2,1,\ldots)$, the exact Riemannian second variation is
\begin{equation}
 H_{n+s}=Q^T\left[H_{\rm emb}
 -\diag(\lambda_1I_4,0,\lambda_2I_4,0,\ldots)\right]Q.
 \label{eq:riemannian-hessian}
\end{equation}
For artificial control (c), rows and columns belonging to the 81 fixed core
tangent variables are deleted. This restriction is an exact Hessian of
Eq.~\eqref{eq:lattice-energy} on that constrained manifold. It is not a
full-grid physical Hessian. On every grid, an independent centered geodesic
second difference with all fixed-core tangent entries set to zero reproduces
$v^TH_{n+s}v$; the largest best-step relative residual is
$2.44\times10^{-9}$.

On the identical interior Dirichlet grid, define
\begin{align}
 H_\Delta&=I_2\otimes\left[L_D+m_\Delta^2
 +\chi(1-n_0)\right],
 &m_\Delta^2&=0.81,&\chi&=-0.30,\label{eq:diquark}\\
 H_A&=I_3\otimes L_D,& H_{\bar cc}&=L_D,
 &&(\xi=1).\label{eq:gauge-ghost}
\end{align}
These are declared quadratic blocks. Here $H_A$ has three real spatial
components and $H_{\bar cc}$ one complex ghost: they are Abelian Feynman-gauge
proxies \cite{FaddeevPopov1967APE7}. A physical SU(2) regulator would instead
require at least $3\times3M=9M$ real spatial-adjoint gauge variables and $3M$
complex adjoint ghosts, plus interacting cross blocks. Those are absent.

For an explicit same-grid fermion proxy, let
\begin{align}
 D_W={}&\sum_{i=1}^3\gamma_i\frac{T_i^+-T_i^-}{2a}+m\one
 +\frac{r_W}{2a}\sum_{i=1}^3(2-T_i^+-T_i^-)\nonumber\\
 &+y\left(P_R\otimes U_x+P_L\otimes U_x^\dagger\right),
 \qquad U_x=n_{x0}\one+i n_{xa}\tau_a,
 \label{eq:wilson-yukawa}
\end{align}
with $m=0.60$, $y=0.80$, $r_W=1$, and Dirichlet shifts. This is a
Wilson/Yukawa operator in the sense of an explicitly declared regulator
\cite{Wilson1974APE7}. We diagonalize
\begin{equation}
 K_F=D_W^\dagger D_W\succeq0.
 \label{eq:normal}
\end{equation}
The displayed same-grid direct sum is therefore
\begin{equation}
 \cH_{\rm declared}=
 H_{n+s}^{\rm core\,restricted}\oplus H_\Delta\oplus H_A
 \oplus H_{\bar cc}\oplus K_F.
 \label{eq:direct-sum}
\end{equation}
Only the first block is a bosonic Hessian of the interacting background
energy. $H_\Delta$ is the Hessian of a separately declared quadratic sector;
$H_A$ is a free gauge proxy; $H_{\bar cc}$ is a Grassmann operator; and $K_F$
is a positive normal proxy. Equation~\eqref{eq:direct-sum} is not the BRST
super-Hessian nor the second variation of $-\log\det D_W$.

\section{Numerical results}

The deterministic card records five independently re-relaxed grid/volume
pairs:
\begin{equation}
 (N,L)=(7,6),(9,6),(11,6),(9,8),(11,10).
 \label{eq:grid-sequence}
\end{equation}
The first three vary $a=L/(N-1)$ at fixed volume; $(7,6)$, $(9,8)$, and
$(11,10)$ give the fixed-$a=1$ volume sequence. No result is obtained by
interpolating a previously relaxed lattice field.

\subsection{Unconstrained unwind}

Table~\ref{tab:unwind} shows that full-grid solves converge to the exact
vacuum within numerical tolerance. Here $g_a$ is the sitewise tangent-gradient
RMS divided by $a^3$.

\begin{table}[htbp]
\centering
\small
\caption{Independent full-grid minimizations of Eq.~\eqref{eq:lattice-energy}.}
\label{tab:unwind}
\begin{tabular}{rrrrrrr}
\toprule
$N$&$L$&$a$&$E_{\rm initial}$&$E_\star$&$g_a$&$B_{\rm geom}$\\
\midrule
7&6&1.00&43.95516354&$9.07\times10^{-16}$&$1.99\times10^{-9}$&0\\
9&6&0.75&54.15076270&$5.33\times10^{-17}$&$3.91\times10^{-9}$&0\\
11&6&0.60&61.09804327&$8.07\times10^{-16}$&$3.17\times10^{-9}$&0\\
9&8&1.00&43.16126382&$1.77\times10^{-14}$&$1.84\times10^{-8}$&0\\
11&10&1.00&42.97909652&$2.04\times10^{-15}$&$2.13\times10^{-9}$&0\\
\bottomrule
\end{tabular}
\end{table}

Fixing only the central value $n(0)=-1$ also fails. It converges to
$E=12.46071222$, $g_a=7.15\times10^{-9}$, but all three regular targets give
$B_{\rm geom}=0$ and $\delta_a=0$. The single pinned point is surrounded by a
degenerate tetrahedral dislocation.

\subsection{Admissible-component guard}

At $N=7,L=6$, every accepted guard endpoint has three-target $B=1$, yet each
search terminates on its floor with nonzero gradient
(Table~\ref{tab:guard}). The
monotone fall of $E$ and $g_a$ as the floor is lowered is evidence that the
descent is driven toward the dislocation boundary, not toward a resolved
interior critical point. This is a controlled numerical result, not a proof
that no other metastable critical point exists anywhere in the component.

\begin{table}[htbp]
\centering
\small
\caption{Admissible-component descent at $N=7,L=6$.}
\label{tab:guard}
\begin{tabular}{rrrrrr}
\toprule
$\delta_{\rm floor}$&steps&$E_\star$&$g_a$&final $\delta_a$&$B_{\rm geom}$\\
\midrule
0.10&21&28.63959591&0.47292019&0.100000002&1\\
0.05&25&26.04401742&0.43767430&0.050000002&1\\
0.02&28&24.72976823&0.41784577&0.020000002&1\\
\bottomrule
\end{tabular}
\end{table}

\subsection{Artificial fixed-core control}

The free residual passes the $2\times10^{-6}$ tolerance on every artificial
fixed-core grid, while the full unanchored residual does not
(Table~\ref{tab:anchor-spectrum}). This separation is why the result is never
promoted to full stationarity.

\begin{table}[htbp]
\centering
\scriptsize
\caption{Artificial controls and lowest same-grid spectra.}
\label{tab:anchor-spectrum}
\begin{tabular}{rrrrrrrrr}
\toprule
$N$&$L$&free $g_a$&full $g_a$&$B_{\rm FD}$&
$\lambda_0^{n+s}$&$\lambda_0^\Delta$&$\lambda_0^A$&$\sigma_{\min}D_W$\\
\midrule
7&6&$1.45\times10^{-8}$&0.5792&0.16647&2.26434&1.52463&0.80385&1.50522\\
9&6&$3.33\times10^{-8}$&0.6631&0.34779&1.80433&1.53862&0.81195&1.46312\\
11&6&$6.11\times10^{-8}$&0.6379&0.47754&1.50021&1.55084&0.81572&1.43240\\
9&8&$9.46\times10^{-9}$&0.3453&0.16638&1.16107&1.21935&0.45672&1.44662\\
11&10&$1.02\times10^{-8}$&0.2362&0.16644&0.75628&1.07593&0.29366&1.42132\\
\bottomrule
\end{tabular}
\end{table}

The largest exact restricted $n+s$ matrix occurs at $N=11$: it has 2,835
degrees of freedom and 77,031 nonzeros. On that grid the real diquark block
has 1,458 degrees of freedom and 9,234 nonzeros, the gauge block has 2,187 and
13,851, the complex ghost has 729 and 4,617, and the fermion normal operator
has 5,832 complex degrees of freedom and 190,432--190,636 nonzeros over the
two $N=11$ backgrounds.

The $N=7,L=6$ artificial solve was independently repeated; canonical JSON
serialization gave the identical digest
\begin{center}
\path{27c5e4f06d5a27dc8866bd8bad32c8205a0d5f9d7ada10a9eb1dfecb6c4ca6c8}.
\end{center}
The complete machine-readable values, source hashes, and guard traces are in
\path{route_f/output/ap_e7_discrete_rerelax_superhessian.json}.

\begin{gate}[What the fixed-core spectrum means]
Positive eigenvalues of $H_{n+s}^{\rm core\,restricted}$ establish only local
stability under variations that vanish on the artificial core and outer
boundary. They do not establish a full-grid $B=1$ stationary point. The smooth
derivative estimator $B_{\rm FD}$ is also not integer on these coarse grids;
it is a convergence diagnostic, never a replacement for admissible geometric
degree.
\end{gate}

\section{Fail-closed conclusion and next constructions}

The finite-site theorem and numerical unwind close an ambiguity left by
AP-E6: simply re-running a generic optimizer cannot produce a genuine
topological lattice soliton for the declared naive action. At least one new
ingredient is necessary. Three controlled options are:
\begin{enumerate}
\item add an infinite or systematically removable barrier
  $V_{\rm adm}(\delta_T)$ and prove observables become barrier-independent;
\item replace the naive Skyrme cell by a topology-preserving spherical-volume
  action and repeat the $a,L$ scan;
\item retain the continuum-relaxed soliton, but discretize only its fluctuation
  operator with a convergence theorem that never promotes finite-site degree
  to an exact sector.
\end{enumerate}
Only after one construction produces an unanchored stationary sequence may
the fermion contribution be upgraded from $D_W^\dagger D_W$ to the bosonic
second derivative of the regulated determinant and combined with interacting
SU(2) gauge--ghost cross blocks.

Accordingly the genuine relaxed-$B=1$, full unanchored stationarity, physical
aggregate super-Hessian, four-dimensional dynamics, HMC, quantum continuum,
physics-promotion, and portal-start flags are all false.

\bibliographystyle{unsrt}
\bibliography{ap_e7_discrete_rerelax_superhessian}

\end{document}
